\subsection{坐标图、图册与流形}

5.1.1. 设 X 是一个集合。X 上的一个坐标图是三元组 \(c = (U, \varphi, E)\)，其中 U 是 X 的子集，E 是巴拿赫空间，\(\varphi\) 是从 U 到 E 的开子集的双射。称 U 为坐标图 c 的定义域。若 E 的维数为 n，则称 c 的维数为 n。否则置 \(\dim c = +\infty\)。

5.1.2. 当不会对 r 产生歧义时，若 X 的两个坐标图 \(c = (\mathbf{U}, \varphi, \mathbf{E})\) 和 \(c' = (\mathbf{U}', \varphi', \mathbf{E}')\)\(X\) 满足下列条件，则称它们 \(C^r\)\(r\)-相容（或简写为相容）：

\begin{enumerate}
\def\labelenumi{(\alph{enumi})}
\item
  \(\varphi (\mathbf{U}\cap \mathbf{U}^{\prime})\)（分别地，\(\varphi^{\prime}(\mathbf{U}\cap \mathbf{U}^{\prime}))\)）在 \(\mathbf{E}\)（分别地，\(\mathbf{E}'\)）中是开集；
\item
  从 \(\varphi \circ \varphi'^{-1}\) 到 \(\varphi' \circ \varphi^{-1}\) 的映射 \(\varphi'(U \cap U')\)（分别地，从 \(\varphi(U \cap U')\) 到 \(\varphi(U \cap U')\) 的映射 \(\varphi'(U \cap U')\)）属于 \(C^r\) 类（参见 2.3.1、3.2.1 和 4.2.1）。
\end{enumerate}

5.1.3. 称集合 X 的一个 \(\mathbf{C}^{r}\)-图册（或简写为图册）为 X 的一族坐标图，若这些坐标图两两 \(\mathbf{C}^{r}\)-相容且其定义域的并为 X。若 \(\mathbf{C}^{r}\) 是图册，则称 X 的两个图册 A 和 B \(A \cup B\)-等价。图册之间的 \(\mathbf{C}^{r}\)-等价关系是等价关系。

5.1.4. 设 \(\mathfrak{S}\) 是巴拿赫空间的集合。若图册 \(\mathcal{A}\) 的每个坐标图 c = (U, , E) 都满足 E  \(\mathfrak{S}\)，则称图册 \(E \in \mathfrak{S}\) 为 \(c = (U, \varphi, E)\) 型。同样，若图册 \(\mathcal{A}\) 的每个坐标图 \(\mathcal{A}\) 中的 \(E\)\((U, \varphi, E)\) 都是希尔伯特空间（分别地，具有可数型的希尔伯特空间），则称 \(\mathcal{A}\) 为希尔伯特型（分别地，具有可数维的希尔伯特型）。

5.1.5. 称一个 \(\mathbf{C}^{r}\) 类 K-流形（或 K 上的 \(\mathbf{C}^{r}\) 类流形；当 K 和 r 不会引起歧义时也简写为流形）为赋予了关于 \(\mathbf{C}^{r}\)-等价关系的一类等价图册的集合 X（Sets，第 II 章，第 6 节，第 9 号）。此类中的图册称为流形 X 的图册。属于 X 的一个图册的坐标图称为流形 X 的坐标图。定义域包含点 a \(\in X\) 的 X 的坐标图称为 X 在 a 处的坐标图。以 a 为中心的坐标图是 X 在 a 处的坐标图 \(\varphi\)，满足 \(\varphi(a) = 0\)。

若 X 是集合，A 是 X 的图册，则赋予 A 的等价类的集合 X 称为由 A 定义的流形。

当 \(r \neq \omega\)（因而 \(K = R\)）时，\(C^r\) 类流形有时称为可微流形。\(C^\omega\) 类流形也称为 \(K\) 上的解析流形（或 \(K\)-解析流形）。

此外，当 K = R（分别地，C、\(Q_{p}\)）时，也称实解析流形（分别地，复解析流形、p-进解析流形）。

5.1.6. 设 X 是流形。X 中作为 X 的坐标图定义域之并的子集构成 X 上某个拓扑的开集族，该拓扑称为流形结构的底层拓扑。对于 X 的每个坐标图 \(c = (U, \varphi, E)\)，映射 \(\varphi\) 是从赋予 X 所诱导拓扑的开集 U 到 E 的开集 \(\varphi(U)\) 的同胚。

X 的底层拓扑空间是贝尔空间。当 K 等于 R 或 C 时，它是局部连通的。

设 X 是流形，A 是 X 的图册。拓扑空间 X 为豪斯多夫空间的充要条件是满足下列条件：对于图册 A 中任意坐标图 (U，\(\varphi\)，E) 和 (V，\(\psi\)，F)，从 \(\psi \circ \varphi^{-1}\) 到 \(\varphi(\mathrm{U} \cap \mathrm{V})\) 的映射 \(\psi(\mathrm{U} \cap \mathrm{V})\) 的图像在 \(\varphi(\mathrm{U}) \times \psi(\mathrm{V})\) 中是闭的。

设 X 是流形；假设拓扑空间 X 是正则的。则对每个点 \(a \in X\)，存在 X 在 a 处的一个坐标图 (U，\(\varphi\)，E)，具有如下性质：U 的子集 Y 在 X 中闭的充要条件是其像 \(\varphi(Y)\) 在 E 中闭。若空间 X 是仿紧的，则 X 上存在一个与 X 的拓扑相容的距离，并使 X 成为完备度量空间。

5.1.7. 设 S 是巴拿赫空间的集合。若一个流形具有 S 型（分别地，希尔伯特型、可数维希尔伯特型）图册，则称其为 S 型流形（分别地，希尔伯特型流形、可数维希尔伯特型流形）。若 S 约化为单个元素 E，则 S 型流形也称为 E 型纯流形。

维数为 n 的纯流形称为 \(K^{n}\) 型纯流形。若流形是 S 型的，其中 \(S = \{K^{n}; n \in N\}\)，则称其局部有限维。

5.1.8. 设 X 是流形，且 \(x \in X\)。X 在 x 处的坐标图的维数（有限或 \(+\infty\)）只取决于 x。称其为 X 在 x 处的维数，记作 \(\dim_{x}X\)。X 的维数记作 \(\dim X\)，是 \(\dim_{x}X\) 时 \(x \in X\) 的上界。

函数 \(x \mapsto \dim_x X\) 局部为常值。流形 \(X\) 局部有限维（分别地，纯 \(n\) 维）的充要条件是对所有 \(\dim_x X < +\infty\) 都有 \(\dim_x X = n\)（分别地，\(x \in X\)）。

5.1.9. 假设 K 局部紧。设 X 是分离流形。X 局部有限维的充要条件是 X 局部紧。

5.1.10. 设 X 是流形，\(\xi^{1},\ldots,\xi^{n}\) 是从 X 的子集 U 到 K 的映射。若三元组 (U，\(\xi=(\xi^{1},\ldots,\xi^{n})\)，K \(\xi\)) 是 X 的坐标图，则称 \(^{n}\) 是 X 在 U 中的坐标系；该坐标图也记作 (U；\(\xi\)) 或 (U；\(\xi^{1},\ldots,\xi^{n}\))。若 \(a\in\mathbf{U}\)，也称 \(\xi\) 是 X 在 a 处的坐标系；若进一步对 \(\xi^{i}(a)=0\) 有 \(i=1,2,\ldots,n\)，则称坐标系 \(\xi\) 以 a 为中心。

